% TikZ 图：验证动量守恒实验
\begin{tikzpicture}[>=Latex,scale=0.96,every node/.style={font=\scriptsize}]
  % Before collision.
  \begin{scope}
    \node[font=\small\bfseries] at (0.42,1.32) {碰前};
    \draw[thick] (0,0) -- (4.90,0);
    \draw[gray!60,thick] (0.15,-0.08) -- (4.75,-0.08);

    \draw[thick,fill=blue!10] (0.65,0.10) rectangle (1.55,0.62);
    \node at (1.10,0.36) {$m_1$};
    \draw[->,red,thick] (1.68,0.42) -- (2.65,0.42)
      node[midway,above] {$v_1$};

    \draw[thick,fill=orange!12] (3.25,0.10) rectangle (4.15,0.62);
    \node at (3.70,0.36) {$m_2$};
    \node[above] at (3.70,0.72) {$v_2$ 可为 $0$};

    \draw[dashed,thick] (2.25,-0.05) -- (2.25,1.08);
    \node[above] at (2.25,1.08) {光电门};
  \end{scope}

  % After collision.
  \begin{scope}[xshift=6.10cm]
    \node[font=\small\bfseries] at (0.42,1.32) {碰后};
    \draw[thick] (0,0) -- (4.90,0);
    \draw[gray!60,thick] (0.15,-0.08) -- (4.75,-0.08);

    \draw[thick,fill=blue!10] (0.62,0.10) rectangle (1.52,0.62);
    \node at (1.07,0.36) {$m_1$};
    \draw[->,blue,thick] (1.65,0.26) -- (2.42,0.26)
      node[midway,below] {$v_1'$};

    \draw[thick,fill=orange!12] (2.82,0.10) rectangle (3.72,0.62);
    \node at (3.27,0.36) {$m_2$};
    \draw[->,orange!80!black,thick] (3.85,0.52) -- (4.65,0.52)
      node[midway,above] {$v_2'$};

    \draw[dashed,thick] (2.25,-0.05) -- (2.25,1.08);
    \draw[dashed,thick] (4.25,-0.05) -- (4.25,1.08);
    \node[above] at (3.25,1.08) {光电门记录碰后速度};
  \end{scope}

  \node[font=\small,align=center] at (5.50,-0.72)
    {统一规定正方向，速度代数值带符号；比较碰前、碰后系统总动量};
\end{tikzpicture}
